\section{An exponential pulse-broadening function implies $\alpha=4$}
\label{app:emg-alpha4}

Exponentially modified Gaussian (EMG) profiles---a Gaussian pulse of width
$\sigma$ convolved with the one-sided exponential
$p(t)=\tau^{-1}e^{-t/\tau}$ ($t\ge0$)---are the standard phenomenological
scattering model in the pulsar and FRB literature, and they are the kernel our
pipeline evaluates in its exponential limit,
\begin{equation}
\label{eq:emg}
I(t)\;\propto\;\frac{1}{2\tau}\,
\exp\!\left[\frac{\sigma^{2}}{2\tau^{2}}-\frac{t-t_{0}}{\tau}\right]
\operatorname{erfc}\!\left[\frac{1}{\sqrt{2}}
\left(\frac{\sigma}{\tau}-\frac{t-t_{0}}{\sigma}\right)\right].
\end{equation}
This appendix records why adopting an EMG kernel is not scaling-agnostic:
within thin-screen scattering the exponential pulse-broadening function (PBF)
is uniquely the $\beta=4$ member of the power-law family of
Section~\ref{sec:beta-scattering-methods}, and $\beta=4$ fixes $\alpha=4$.
Fixing $\alpha=4$ in an EMG fit is therefore a self-consistency requirement
rather than a modeling choice, and conversely an EMG fit with $\alpha$ free
combines incompatible turbulence assumptions.

\subsection{Uniqueness of the exponential limit}

For a thin screen with $P_n(q)\propto q^{-\beta}$, $2<\beta<4$, and a
negligible inner scale, the PBF is exponential at small delays and develops a
power-law tail at large delays \citep{Cordes2025}:
\begin{equation}
\label{eq:pbf-family}
p(s)\;\propto\;
\begin{cases}
e^{-s}, & s\le s_{c},\\[2pt]
e^{-s_{c}}\,(s/s_{c})^{-\beta/2}, & s>s_{c},
\end{cases}
\qquad
s_{c}=2\ln\!\frac{2}{4-\beta},
\end{equation}
with $s=t/\tau$. The fraction of the unit-normalized PBF carried by the tail
is
\begin{equation}
\label{eq:tailmass}
M(\beta)\;=\;
\int_{s_{c}}^{\infty} e^{-s_{c}}\!\left(\frac{s}{s_{c}}\right)^{-\beta/2} ds
\;=\;\frac{(4-\beta)^{2}\,s_{c}}{2\,(\beta-2)},
\end{equation}
which is strictly positive for every $2<\beta<4$ (e.g.\ $M(3)=\ln 2$) and
vanishes only in the limit $\beta\to4$, where $s_{c}\to\infty$ and the family
degenerates to the pure exponential. The tail is normalizable for $\beta>2$
while its mean delay converges only for $\beta>4$, so every $\beta<4$ member is
qualitatively heavier-tailed than the exponential rather than a small
perturbation of it. An exactly exponential PBF therefore corresponds to one and
only one point of the family: the square-law limit $\beta=4$.

\subsection{The scaling forced by $\beta=4$}

The same $\beta$ sets the frequency scaling. With the inertial-range phase
structure function $D_\phi(\rho)\propto\nu^{-2}\rho^{\beta-2}$
(Section~\ref{sec:beta-scattering-methods}), the diffractive scale defined by
$D_\phi(\rho_{\rm d})=1$ obeys $\rho_{\rm d}\propto\nu^{2/(\beta-2)}$, the
scattering angle scales as $\theta_{\rm d}\propto(k\rho_{\rm d})^{-1}\propto
\nu^{-\beta/(\beta-2)}$, and the geometric delay $\tau\propto\theta_{\rm d}^{2}$
gives
\begin{equation}
\label{eq:alpha-closure-appendix}
\tau\;\propto\;\nu^{-2\beta/(\beta-2)}
\quad\Longrightarrow\quad
\alpha(\beta)=\frac{2\beta}{\beta-2},
\qquad \alpha(4)=4 .
\end{equation}
An EMG kernel thus commits a fit to $\tau\propto\nu^{-4}$ by definition.

Three corollaries govern our fit design. First, the only fixed scattering index
consistent with an exponential kernel is $\alpha=4$; fixing the Kolmogorov
value $\alpha=22/5$ inside an EMG would be inconsistent, because $\beta=11/3$
implies the tailed PBF of Equation~\ref{eq:pbf-family}, not an exponential.
Second, an EMG fit with $\alpha$ free asserts $\beta=4$ through its shape while
asserting $\beta\neq4$ through its scaling; a shallow apparent index recovered
that way measures departure from the assumed scattering family---burst
morphology, an extended medium, or a finite inner scale
\citep{Bhat2004,Cordes2025}---rather than a turbulence spectral index. Third,
$\alpha=4$ is a two-sided \emph{floor} for the whole single-thin-screen family:
Equation~\ref{eq:alpha-closure-appendix} exceeds 4 for every $\beta<4$
($\alpha$ is monotonically decreasing in $\beta$), while for $\beta\ge4$ the
structure function saturates to the square-law form and pins
$\tau\propto\nu^{-4}$ regardless of $\beta$ \citep{Rickett1990}. An empirical
$\alpha<4$ is therefore unreachable within this family and must be read as a
violated assumption, never as a spectral index
(Section~\ref{sec:beta-scattering-methods}). This is
why our primary model samples $\beta$ directly and derives both the PBF shape
and $\alpha(\beta)$ from it (Section~\ref{sec:beta-scattering-methods}).

% The algebra in this appendix (closure values, limiting behavior, tail
% convergence conditions, and equivalences of the EMG forms) was verifiedf
% symbolically with WolframScript.
